\chapter{Listing a directory: \cmd{ls} from the inode up}
\label{ch:ls}

You have typed \cmd{ls} thousands of times. By the end of this chapter you will know what it asks the kernel for, why \cmd{ls -l} is slower than plain \cmd{ls}, and why it hangs on a directory with two million files.

\begin{origin}[1971, Unix V1]
The first Unix manual already listed \cmd{ls}. Its job was small: read a directory and print the names inside it. Everything it has grown since then is an answer to one question: \emph{what else might you want to know about those names?}
\end{origin}

\section{A directory is a file of names}

\begin{plainwords}
A directory is a list. Each line of the list holds a name and a number. The number points to the real file.
\end{plainwords}

\begin{definition}{Directory entry (dentry)}
A record inside a directory that pairs a file name with an \term{inode number}. The inode, not the name, holds the file's metadata and the location of its data.
\end{definition}

\Cref{fig:dentry-inode} shows the three layers. Plain \cmd{ls} only needs the left box; \cmd{ls -l} must visit the middle box once per entry.

\svgfigure{ch01-dentry-inode}{A directory maps names to inode numbers; the inode holds metadata and points to data blocks. Notice which layer each command has to touch.}{fig:dentry-inode}

\begin{underthehood}
Run \cmd{strace -e trace=getdents64,statx ls -l} on a small directory. You will see one \cmd{getdents64} call that returns all the names, then one \cmd{statx} call per file. That second loop is the cost of \flag{-l}.
\end{underthehood}

\begin{shell}
$ strace -c -e trace=getdents64,statx ls -l /etc > /dev/null
% time     seconds  usecs/call     calls    errors syscall
------ ----------- ----------- --------- --------- ----------------
 91.40    0.000851           3       241           statx
  8.60    0.000080          27         3           getdents64
\end{shell}

\section{The flags, grouped by what they change}

\Cref{tab:ls-flags} groups the flags you will meet most. The full list is in the GNU manual~\cite{gnu-ls}.

\begin{table}[htbp]
\centering
\caption{Common \cmd{ls} flags, grouped by purpose (GNU coreutils).}
\label{tab:ls-flags}
\begin{tabularx}{\linewidth}{@{}l l X@{}}
\toprule
Group & Flag & What it does \\
\midrule
What to show & \flag{-a} & Include entries starting with a dot, including \cmd{.} and \cmd{..} \\
             & \flag{-A} & Like \flag{-a} but leaves out \cmd{.} and \cmd{..} \\
Detail       & \flag{-l} & Long format: type, mode, links, owner, size, time \\
             & \flag{-i} & Print each entry's inode number \\
Sizes        & \flag{-h} & Human-readable sizes, powers of 1024 \\
             & \flag{--si} & Human-readable sizes, powers of 1000 \\
Sorting      & \flag{-t} & Newest first, by modification time \\
             & \flag{-U} & Do not sort: directory order, fastest on huge directories \\
\bottomrule
\end{tabularx}
\end{table}

\begin{gotcha}[Sorting needs everything first]
\cmd{ls} cannot print the first name until it has read every name, because sorting needs the whole list. On a directory with millions of entries that looks like a hang. \flag{-U} skips the sort and streams.
\end{gotcha}

\begin{interview}
``Why is \cmd{ls -l} slower than \cmd{ls} on a network filesystem?'' A strong answer names the extra \cmd{statx} per entry and notes that each one may be a network round-trip.
\end{interview}

\chaptersummary
\begin{itemize}
  \item A directory stores names and inode numbers; the inode stores everything else.
  \item \flag{-l} costs one \cmd{statx} per entry; sorting costs reading every entry first.
\end{itemize}

\chaptervocab
\begin{vocablist}
  \item[Inode] The record holding a file's metadata and block locations.
  \item[Directory entry] A name-to-inode pair stored inside a directory.
\end{vocablist}

\chapterexercises
\begin{exercise}[1]{ex:ls-a-vs-A}
What is the difference between \flag{-a} and \flag{-A}?
\end{exercise}
\begin{exercise}[3]{ex:ls-hang}
A colleague says \cmd{ls} ``hangs'' in \file{/var/spool/queue}. Explain why in terms of system calls, and give two commands that would list the directory without hanging.
\end{exercise}

\furtherreading
\begin{itemize}
  \item The GNU \cmd{ls} manual~\cite{gnu-ls}: read the sorting section closely.
\end{itemize}
